%% GENERATED by tools/gen_latex.py -- edit content/ch09.py instead.
%% Build:  latexmk -pdf ch09.tex     (or: pdflatex ch09.tex, twice)
\documentclass[aspectratio=169,11pt,t]{beamer}
\usepackage{itbpro}

\coursetitle{ConvNet Architecture Patterns}
\coursesubtitle{Four best practices that turn a working image model into a good one --- and a mini Xception that beats the previous chapter's model with less than half the parameters.}
\coursekicker{Chapter 9}
\coursesource{Chollet \& Watson, Deep Learning with Python 3e --- chapter 9}
\courseduration{2.5 hours}
\coursepresenter{Rahman Indra Kesuma, S.Kom., M.Cs. $\cdot$ Prof. Bambang Riyanto Trilaksono}{Teaching Assistant $\cdot$ Lead Instructor}
\coursebrand{AI for Professional \textperiodcentered\ ITB \textperiodcentered\ Ch 9}

\title{ConvNet Architecture Patterns}
\author{Rahman Indra Kesuma, S.Kom., M.Cs. $\cdot$ Prof. Bambang Riyanto Trilaksono}
\date{}

\begin{document}
\covertitle

\framekicker{Learning outcomes}
\begin{frame}[shrink=25]{Session Objectives}
\leadpar{By the end of this session, participants can:}
\begin{isteps}
\item Explain architecture as a \textbf{choice of hypothesis space}, and why a good one reduces the search gradient descent has to do.
\item Apply the \textbf{modularity--hierarchy--reuse} formula to read any published ConvNet.
\item Explain the \textbf{vanishing gradient} problem and fix it with \textbf{residual connections}, including the 1$\times$1 projection when shapes differ.
\item Say what \textbf{batch normalisation} does, where to put it, and why \inlinecode{use\_bias=False} accompanies it.
\item Explain the assumption behind \textbf{depthwise separable convolutions}, and why they are not faster on a GPU despite using far fewer parameters.
\item Assemble all four into a \textbf{mini Xception}, and read the resulting parameter-count and accuracy change.
\item Say when a \textbf{Vision Transformer} is the better choice, and when it is not.
\end{isteps}

\vskip 6pt
\vskip 4pt
\reslink{Book}{Chapter 9 --- full text}{https://deeplearningwithpython.io/chapters/chapter09\_convnet-architecture-patterns}
\reslink{Notebook}{01\_residual\_connections.ipynb}{https://hendrikarisma.my.id/rs/ai-products-course/notebooks/ch09/01\_residual\_connections.html}
\reslink{Notebook}{02\_batchnorm\_and\_separable.ipynb}{https://hendrikarisma.my.id/rs/ai-products-course/notebooks/ch09/02\_batchnorm\_and\_separable.html}
\reslink{Notebook}{03\_mini\_xception.ipynb}{https://hendrikarisma.my.id/rs/ai-products-course/notebooks/ch09/03\_mini\_xception.html}
\reslink{Notebook}{All chapter 9 notebooks}{https://hendrikarisma.my.id/rs/ai-products-course/notebooks/index.html\#ch09}
\reslink{Repo}{Author's official notebook}{https://github.com/fchollet/deep-learning-with-python-notebooks/blob/master/chapter09\_convnet-architecture-patterns.ipynb}
\vskip 2pt
\end{frame}

\framekicker{Introduction}
\begin{frame}[shrink=25]{Architecture is the choice of hypothesis space}
{\fontsize{9.4}{12.6}\selectfont\color{fg2}A model's architecture is the sum of the choices that went into it: which layers, configured how, connected in what arrangement. Those choices \textbf{define the hypothesis space} --- the set of functions gradient descent may search.\par}\vskip 4pt
\begin{iband}[signal]
Like feature engineering, a good hypothesis space \textbf{encodes prior knowledge}. Using convolution layers means you already know that the relevant patterns are \hilite{translation invariant}
\end{iband}
\end{frame}

\framekicker{Introduction}
\begin{frame}[shrink=25]{Why it is often the difference between success and failure}
\begin{minipage}[t]{0.494\textwidth}
\begin{icardbad}
\cardh{Bad choices cannot be trained away}
The model gets stuck at suboptimal metrics, and \textbf{no amount of training data will save it}.
\end{icardbad}
\end{minipage}%
\hfill
\begin{minipage}[t]{0.494\textwidth}
\begin{icardgood}
\cardh{Good choices multiply your data}
Learning accelerates and the model uses the available data efficiently, \textbf{reducing the need for a large dataset}.
\end{icardgood}
\end{minipage}%
\begin{iquote}
\quotetext{A good model architecture is one that reduces the size of the search space, or otherwise makes it easier to converge to a good point of the search space.}\quotecite{Chollet \& Watson, chapter 9 introduction}
\end{iquote}
\end{frame}

\framekicker{Introduction}
\begin{frame}[shrink=25]{An honest note about how this knowledge is held}
{\fontsize{9.4}{12.6}\selectfont\color{fg2}Model architecture is \textbf{more an art than a science}. Experienced engineers assemble high-performing models on the first try; beginners struggle to build one that trains at all.\par}\vskip 4pt
\begin{iband}[amber]
The keyword is \textit{intuitively}: \textbf{no one can give you a clear explanation of what works and what does not}. Experts rely on pattern matching acquired through practice.
\end{iband}
{\fontsize{9.4}{12.6}\selectfont\color{fg2}But it is not \textit{only} intuition. As in any engineering discipline there are \textbf{best practices}, and this chapter is four of them. Remember, too, that gradient descent is \hilite{a pretty stupid search process} --- it needs all the help it can get.\par}\vskip 4pt
\end{frame}

\coursesection{01}{Modularity, hierarchy, and reuse}{The universal recipe for making a complex system simpler.}

\framekicker{Section 9.1}
\begin{frame}[shrink=25]{The MHR formula}
\begin{center}
\includegraphics[width=0.92\linewidth,height=0.60\textheight,keepaspectratio]{figs/ch09-mhr.tex.pdf}
\end{center}
\vskip 3pt{\color{ink3}\fontsize{7.4}{9.4}\selectfont Figure 9.1 --- the same recipe underlies a cathedral, your body, a navy, and the Keras codebase.\par}
{\fontsize{9.4}{12.6}\selectfont\color{fg2}Software engineers already know this: an effective codebase is modular, hierarchical, and does not reimplement the same thing twice. Following those principles \textbf{is} software architecture.\par}\vskip 4pt
\end{frame}

\framekicker{Section 9.1}
\begin{frame}[shrink=25]{Deep learning is that recipe applied to optimisation}
\begin{ibullets}
\item Take a classic optimisation technique --- \textbf{gradient descent over a continuous function space}.
\item Structure the search space into \textbf{modules}: layers.
\item Organise them into a \textbf{deep hierarchy} --- often just a stack, the simplest kind.
\item \textbf{Reuse} whatever you can: convolutions are entirely about reusing the same information at \hilite{different spatial locations}.
\end{ibullets}
\end{frame}

\framekicker{Section 9.1}
\begin{frame}[shrink=25]{Two patterns visible in every published ConvNet}
\begin{minipage}[t]{0.494\textwidth}
\begin{icardaccent}
\cardh{Repeated blocks}
Popular architectures are structured not merely into layers but into \textbf{repeated groups of layers} --- \textit{blocks} or \textit{modules}. Xception repeats a SeparableConv -- SeparableConv -- MaxPooling block.
\end{icardaccent}
\end{minipage}%
\hfill
\begin{minipage}[t]{0.494\textwidth}
\begin{icardaccent}
\cardh{A feature pyramid}
Filter count \textbf{grows with depth} while feature-map size \textbf{shrinks}. You saw 32 $\rightarrow$ 64 $\rightarrow$ 128 in chapter 8; the same pattern runs through Xception.
\end{icardaccent}
\end{minipage}%
{\fontsize{9.4}{12.6}\selectfont\color{fg2}Once you can see those two patterns, most architecture diagrams in papers become readable at a glance.\par}\vskip 4pt
\end{frame}

\framekicker{Section 9.1 $\cdot$ figure 9.2}
\begin{frame}[shrink=25]{Xception, read with those two patterns}
\begin{center}
\includegraphics[width=0.92\linewidth,height=0.60\textheight,keepaspectratio]{figs/ch09-xception.tex.pdf}
\end{center}
\vskip 3pt{\color{ink3}\fontsize{7.4}{9.4}\selectfont Repeated blocks, growing depth, and a 1x1 projection carrying each residual across the downsampling.\par}
{\fontsize{9.4}{12.6}\selectfont\color{fg2}The spatial sizes fall 299 $\rightarrow$ 149 $\rightarrow$ 74 $\rightarrow$ 37 while the depth climbs 3 $\rightarrow$ 32 $\rightarrow$ 128 $\rightarrow$ 256 $\rightarrow$ 728. Everything else on this diagram is \hilite{the same block, repeated}\par}\vskip 4pt
\end{frame}

\coursesection{02}{Residual connections}{The game of telephone, and how to stop playing it.}

\framekicker{Section 9.2}
\begin{frame}[shrink=25]{Backpropagation as a game of telephone}
{\fontsize{9.4}{12.6}\selectfont\color{fg2}A message whispered from player to player ends up bearing little resemblance to the original. Backpropagation through a deep chain has the same problem.\par}\vskip 4pt
\begin{center}
\includegraphics[width=0.92\linewidth,height=0.60\textheight,keepaspectratio]{figs/ch09-vanish.tex.pdf}
\end{center}
\vskip 3pt{\color{ink3}\fontsize{7.4}{9.4}\selectfont Each successive function introduces noise into the error signal travelling backwards.\par}
\end{frame}

\framekicker{Section 9.2}
\begin{frame}[shrink=25]{The vanishing gradient problem}
{\fontsize{9.4}{12.6}\selectfont\color{fg2}To adjust \inlinecode{f1} you must percolate error information back through \inlinecode{f2}, \inlinecode{f3}, and \inlinecode{f4}. Each step adds noise.\par}\vskip 4pt
\begin{iband}[rose]
If the chain is too deep, that noise \textbf{overwhelms the gradient information} and backpropagation stops working. \hilite{The model will not train at all.}
\end{iband}
{\fontsize{9.4}{12.6}\selectfont\color{fg2}The fix is to force each function in the chain to be \textbf{non-destructive} --- to retain a noiseless version of what came in.\par}\vskip 4pt
\end{frame}

\framekicker{Section 9.2}
\begin{frame}[shrink=25]{A residual connection is an information shortcut}
\begin{center}
\includegraphics[width=0.92\linewidth,height=0.60\textheight,keepaspectratio]{figs/ch09-residual.tex.pdf}
\end{center}
\vskip 3pt{\color{ink3}\fontsize{7.4}{9.4}\selectfont Figure 9.3 --- add the input of a block back to its output.\par}
{\fontsize{9.4}{12.6}\selectfont\color{fg2}The shortcut routes around destructive or noisy blocks --- those containing ReLU activations or dropout --- so gradient information from early layers can propagate noiselessly. Introduced in 2015 with \textbf{ResNet} (He et al., Microsoft).\par}\vskip 4pt
\end{frame}

\framekicker{Section 9.2 $\cdot$ listing 9.1}
\begin{frame}[shrink=25]{Three lines, and one condition}
{\fontsize{9.4}{12.6}\selectfont\color{fg2}In code it is almost trivial. Save the input, run the block, add the two.\par}\vskip 4pt
\icode{python}{listing 9.1 --- a residual connection}{listings/ch09-01.py}
\begin{iband}[amber]
The condition: adding requires \textbf{the block output and the residual to have the same shape}. That fails as soon as the block changes the filter count or pools --- which is \hilite{most real blocks}
\end{iband}
\end{frame}

\framekicker{Section 9.2 $\cdot$ listing 9.2}
\begin{frame}[shrink=25]{Case 1 --- the block changes the filter count}
{\fontsize{9.4}{12.6}\selectfont\color{fg2}Project the residual with a \textbf{1$\times$1 convolution and no activation}. It is a purely linear reshaping of the channel dimension.\par}\vskip 4pt
\icode{python}{listing 9.2 --- projecting the residual}{listings/ch09-02.py}
{\fontsize{9.4}{12.6}\selectfont\color{fg2}\inlinecode{padding="same"} in the block avoids the spatial shrinkage a convolution would otherwise cause, \hilite{keeping the two shapes aligned}\par}\vskip 4pt
\end{frame}

\framekicker{Section 9.2 $\cdot$ listing 9.3}
\begin{frame}[shrink=25]{Case 2 --- the block pools}
{\fontsize{9.4}{12.6}\selectfont\color{fg2}If the block halves the spatial size, the projection must halve it too. That is what \inlinecode{strides=2} on the 1$\times$1 convolution is for.\par}\vskip 4pt
\icode{python}{listing 9.3 --- matching a max pooling layer}{listings/ch09-03.py}
\begin{iband}[signal]
So the rule is a pair: \textbf{\inlinecode{padding="same"} inside the block, \inlinecode{strides} on the projection} to match whatever downsampling the block performs.
\end{iband}
\end{frame}

\framekicker{Section 9.3.2}
\begin{frame}[shrink=25]{The same budget, spent two ways}
\begin{center}
\includegraphics[width=0.97\linewidth,height=0.56\textheight,keepaspectratio]{figs/ch09-depth-width.tex.pdf}
\end{center}
\vskip 3pt{\color{ink3}\fontsize{7.4}{9.4}\selectfont Parameter counts are computed from the layer sizes shown, not chosen to make the point.\par}
\end{frame}

\framekicker{Section 9.2}
\begin{frame}[shrink=25]{Why "deep and narrow beats broad and shallow"}
{\fontsize{9.4}{12.6}\selectfont\color{fg2}Residual connections are what make that principle usable. Without them, depth is limited by how far a gradient can travel before the noise swallows it.\par}\vskip 4pt
\begin{minipage}[t]{0.494\textwidth}
\begin{icardwarn}
\cardh{Broad and shallow}
Few layers, each very wide. Many parameters, but only one or two levels of representation --- closer to the \textit{shallow learning} of chapter 1.
\end{icardwarn}
\end{minipage}%
\hfill
\begin{minipage}[t]{0.494\textwidth}
\begin{icardgood}
\cardh{Deep and narrow}
Many layers, each modest. A \textbf{long hierarchy of representations}, which is what deep learning is for --- and now trainable, thanks to the shortcuts.
\end{icardgood}
\end{minipage}%
\end{frame}

\coursesection{03}{Batch normalization}{Normalising not just the input, but every intermediate activation.}

\framekicker{Section 9.3}
\begin{frame}[shrink=25]{The question that motivates it}
{\fontsize{9.4}{12.6}\selectfont\color{fg2}You already normalise data before feeding it in. But normalisation may be a concern \textbf{after every transformation the network performs}.\par}\vskip 4pt
\icode{python}{the familiar normalisation}{listings/ch09-04.py}
\begin{iband}[signal]
Even if the data entering a \inlinecode{Dense} or \inlinecode{Conv2D} layer has zero mean and unit variance, there is \textbf{no reason to expect the same of what comes out}. So: \hilite{could normalising intermediate activations help?}
\end{iband}
\end{frame}

\framekicker{Section 9.3}
\begin{frame}[shrink=25]{What BatchNormalization does}
\begin{itable}{>{\raggedright\arraybackslash}p{0.2068\textwidth}>{\raggedright\arraybackslash}p{0.7332\textwidth}}
\thead{Phase} & \thead{What it normalises with} \\ \midrule
\textbf{Training} & The mean and variance of \textbf{the current batch}. \\
\textbf{Inference} & An \textbf{exponential moving average} of the batchwise mean and variance seen during training --- because a big enough representative batch may not be available at inference time. \\
\end{itable}
{\fontsize{9.4}{12.6}\selectfont\color{fg2}That moving average is exactly the \textbf{non-trainable weight} chapter 7 mentioned: \inlinecode{BatchNormalization} is the only built-in Keras layer that has one.\par}\vskip 4pt
\end{frame}

\framekicker{Section 9.3}
\begin{frame}[shrink=25]{Nobody is quite sure why it works}
{\fontsize{9.4}{12.6}\selectfont\color{fg2}Ioffe and Szegedy's 2015 paper suggested it works by \textit{reducing internal covariate shift}. \textbf{No one really knows for sure.} There are various hypotheses and no certainties.\par}\vskip 4pt
\begin{iquote}
\quotetext{You'll find that this is true of many things in deep learning --- deep learning is not an exact science but a set of ever-changing, empirically derived engineering best practices.}\quotecite{Chollet \& Watson, section 9.3}
\end{iquote}
\begin{iband}[amber]
Worth saying aloud in a classroom: \textbf{"we use it because it measurably helps"} is a legitimate engineering answer, and it is \hilite{the honest one here}
\end{iband}
\end{frame}

\framekicker{Section 9.3}
\begin{frame}[shrink=25]{Where to put it, and one detail that follows}
{\fontsize{9.4}{12.6}\selectfont\color{fg2}The convention is convolution $\rightarrow$ batch norm $\rightarrow$ activation, with the bias term switched off.\par}\vskip 4pt
\icode{python}{the usual ordering}{listings/ch09-05.py}
\begin{iband}[signal]
\inlinecode{use\_bias=False} because batch normalisation \textbf{already centres the output} --- the layer's own bias term would be immediately subtracted away, so it is \hilite{pure waste}
\end{iband}
\end{frame}

\framekicker{Section 9.3}
\begin{frame}[shrink=25]{And why chapter 8 said not to unfreeze it}
\begin{iband}[rose]
During fine-tuning, unfreezing a \inlinecode{BatchNormalization} layer lets its moving statistics be overwritten by a \textbf{small, unrepresentative} batch of your new data --- undoing the calibration the pretrained model arrived at.
\end{iband}
{\fontsize{9.4}{12.6}\selectfont\color{fg2}That is the reason behind the rule you were given in chapter 8 without explanation. \textbf{Leave batch normalisation frozen when fine-tuning.}\par}\vskip 4pt
\end{frame}

\coursesection{04}{Depthwise separable convolutions}{A stronger prior, and far fewer parameters.}

\framekicker{Section 9.4}
\begin{frame}[shrink=25]{The assumption it encodes}
{\fontsize{9.4}{12.6}\selectfont\color{fg2}Regular convolution assumes patterns are \textbf{not tied to specific locations}. Depthwise separable convolution adds a second assumption.\par}\vskip 4pt
\begin{iband}[signal]
That \textbf{spatial locations in intermediate activations are highly correlated, but different channels are highly independent}. That holds generally for the representations deep networks learn, so it serves as \hilite{a useful prior}
\end{iband}
{\fontsize{9.4}{12.6}\selectfont\color{fg2}And the general principle behind it: \textbf{a model with stronger priors about the structure of its information is a better model --- as long as the priors are accurate.}\par}\vskip 4pt
\end{frame}

\framekicker{Section 9.4}
\begin{frame}[shrink=25]{How the operation splits in two}
\begin{center}
\includegraphics[width=0.92\linewidth,height=0.60\textheight,keepaspectratio]{figs/ch09-separable.tex.pdf}
\end{center}
\vskip 3pt{\color{ink3}\fontsize{7.4}{9.4}\selectfont Figure 9.4 --- a depthwise convolution (independent spatial convolutions per channel) followed by a pointwise 1$\times$1 convolution that mixes channels.\par}
\end{frame}

\framekicker{Section 9.4}
\begin{frame}[shrink=25]{The parameter arithmetic}
{\fontsize{9.4}{12.6}\selectfont\color{fg2}Take a 3$\times$3 window with 64 input channels and 64 output channels, and count the trainable parameters both ways.\par}\vskip 4pt
\begin{itable}{>{\raggedright\arraybackslash}p{0.3196\textwidth}>{\raggedright\arraybackslash}p{0.2256\textwidth}>{\raggedright\arraybackslash}p{0.3948\textwidth}}
\thead{Operation} & \thead{Parameter count} & \thead{Arithmetic} \\ \midrule
Regular convolution & \textbf{36,864} & 3 $\times$ 3 $\times$ 64 $\times$ 64 \\
Depthwise separable & \textbf{4,672} & 3 $\times$ 3 $\times$ 64 + 64 $\times$ 64 \\
\end{itable}
\begin{iband}[signal]
Roughly \textbf{eight times fewer}, with comparable representational power --- and the gap \hilite{widens as filters or windows get larger}. Smaller models that converge faster and overfit less.
\end{iband}
\end{frame}

\framekicker{Section 9.4}
\begin{frame}[shrink=25]{So it must be much faster. It is not.}
{\fontsize{9.4}{12.6}\selectfont\color{fg2}On CPU, where the implementation is parallelised C++, you do see a meaningful speedup. On GPU you are not running a simple CUDA implementation --- you are running a \textbf{cuDNN kernel optimised down to each machine instruction}.\par}\vskip 4pt
\begin{iband}[amber]
Despite repeated requests to NVIDIA, depthwise separable convolutions have not received nearly the same optimisation effort, so they remain \textbf{about as fast as regular convolutions} --- while using \hilite{quadratically fewer parameters and operations}
\end{iband}
{\fontsize{9.4}{12.6}\selectfont\color{fg2}Use them anyway: the lower parameter count means less overfitting risk, and the channel-independence assumption gives faster convergence and more robust representations.\par}\vskip 4pt
\end{frame}

\framekicker{Section 9.4}
\begin{frame}[shrink=25]{The wider lesson: hardware, software, and algorithms co-evolve}
{\fontsize{9.4}{12.6}\selectfont\color{fg2}What is a slight inconvenience here becomes an impassable wall elsewhere. The entire ecosystem is micro-optimised for a \textbf{specific set of algorithms} --- ConvNets trained by backpropagation.\par}\vskip 4pt
\begin{iband}[rose]
Experiment with gradient-free optimisation or spiking neural networks and your first CUDA implementations would be \textbf{orders of magnitude slower} than a plain ConvNet, \hilite{no matter how good the idea}. Convincing others to adopt it would be a hard sell even if it were plainly better.
\end{iband}
{\fontsize{9.4}{12.6}\selectfont\color{fg2}GPUs and CUDA enabled backprop-trained ConvNets $\rightarrow$ NVIDIA optimised for them $\rightarrow$ the research community consolidated behind them. \textbf{Taking a different path now would require a multiyear re-engineering of the whole ecosystem.}\par}\vskip 4pt
\end{frame}
\note{A good slide for a discussion about technology strategy: the winning technique is partly the one that got optimised first.}

\coursesection{05}{Putting it together}{A mini Xception, from the four practices.}

\framekicker{Section 9.5}
\begin{frame}[shrink=25]{The six principles, collected}
\begin{isteps}
\item Organise the model into \textbf{repeated blocks} of layers --- usually several convolutions plus a max pooling layer.
\item \textbf{Filter count should increase as feature-map size decreases.}
\item \textbf{Deep and narrow beats broad and shallow.}
\item \textbf{Residual connections} around blocks let you train deeper networks.
\item \textbf{Batch normalisation} after convolution layers is often beneficial.
\item \textbf{\inlinecode{SeparableConv2D} instead of \inlinecode{Conv2D}} is more parameter efficient.
\end{isteps}
\end{frame}

\framekicker{Section 9.5}
\begin{frame}[shrink=25]{One block, with all four practices in it}
\begin{center}
\includegraphics[width=0.92\linewidth,height=0.60\textheight,keepaspectratio]{figs/ch09-block.tex.pdf}
\end{center}
\vskip 3pt{\color{ink3}\fontsize{7.4}{9.4}\selectfont Batch norm, separable convolution, pooling, and a projected residual --- repeated five times with growing depth.\par}
\end{frame}

\framekicker{Section 9.5}
\begin{frame}[shrink=25]{\ldots{}and the whole model is a loop}
{\fontsize{9.4}{12.6}\selectfont\color{fg2}Because the block is uniform, the model is a \inlinecode{for} loop. Note the first layer, and why it is not separable.\par}\vskip 4pt
\icode{python}{the mini Xception}{listings/ch09-06.py}
\begin{iband}[amber]
That first-layer detail is easy to miss and worth stating: \textbf{separable convolutions are wrong for the raw RGB input}, because \hilite{colour channels in natural images are strongly correlated}
\end{iband}
\end{frame}

\framekicker{Section 9.5}
\begin{frame}[shrink=25]{The result: fewer parameters, better accuracy}
\begin{minipage}[t]{0.494\textwidth}
\istat{721,857}{trainable parameters, mini Xception}
\end{minipage}%
\hfill
\begin{minipage}[t]{0.494\textwidth}
\istat{1,569,089}{trainable parameters, chapter 8's model}
\end{minipage}%
\begin{minipage}[t]{0.494\textwidth}
\istat{90.8\%}{test accuracy, mini Xception}
\end{minipage}%
\hfill
\begin{minipage}[t]{0.494\textwidth}
\istat{83.9\%}{test accuracy, chapter 8's model}
\end{minipage}%
\begin{iband}[signal]
\textbf{Less than half the parameters, seven points better.} Following architecture best practices has \hilite{an immediate, sizeable effect}.
\end{iband}
\end{frame}

\framekicker{Section 9.5}
\begin{frame}[shrink=25]{And note what was not done}
{\fontsize{9.4}{12.6}\selectfont\color{fg2}No hyperparameter search was run. The configuration comes \textbf{purely from the best practices above}, plus a small amount of intuition about model size.\par}\vskip 4pt
\begin{iband}[signal]
Systematic hyperparameter tuning is \textbf{chapter 18}. The point here is that \hilite{seven points were available before tuning even started}
\end{iband}
\end{frame}

\framekicker{Section 9.5}
\begin{frame}[shrink=25]{Reusing chapter 8's setup unchanged}
{\fontsize{9.4}{12.6}\selectfont\color{fg2}Only the model definition changed. Data loading, augmentation, callbacks, and the training call are exactly those from section 8.2 --- which is the point of having built them properly once.\par}\vskip 4pt
\icode{python}{the unchanged surroundings}{listings/ch09-07.py}
\begin{iband}[signal]
A useful habit that falls out of this: \textbf{keep the data pipeline and the model definition separable}, so architecture experiments cost \hilite{one edit, not a rewrite}
\end{iband}
\end{frame}

\coursesection{06}{Beyond convolution}{Vision Transformers, and when they are the wrong choice.}

\framekicker{Section 9.6}
\begin{frame}[shrink=25]{What a Vision Transformer does}
{\fontsize{9.4}{12.6}\selectfont\color{fg2}The Transformer was developed to process text --- it is fundamentally a \textbf{sequence-processing} architecture (chapter 15). A ViT applies it to images.\par}\vskip 4pt
\begin{ibullets}
\item It \textbf{splits the image into a 1D sequence of patches}.
\item Each patch becomes a flat vector.
\item The vector sequence is processed like a sentence --- which lets it capture \textbf{long-range relationships} between distant parts of the image, \hilite{something ConvNets can struggle with}.
\end{ibullets}
\end{frame}

\framekicker{Section 9.6}
\begin{frame}[shrink=25]{When to reach for one}
\begin{center}
\includegraphics[width=0.92\linewidth,height=0.60\textheight,keepaspectratio]{figs/ch09-vit.tex.pdf}
\end{center}
\vskip 3pt{\color{ink3}\fontsize{7.4}{9.4}\selectfont The authors' general experience, stated as a decision.\par}
\begin{ibullets}
\item ViTs \textbf{lack the spatial prior} of ConvNets --- the patch-based 2D structure encodes assumptions about local visual structure, which makes ConvNets more data efficient.
\item For ViTs to shine \textbf{they need to be really large}, which makes them unwieldy for anything smaller than ImageNet.
\end{ibullets}
\end{frame}

\framekicker{Section 9.6}
\begin{frame}[shrink=25]{The state of the contest}
\begin{iband}[signal]
\textbf{The battle for image recognition supremacy is far from over}, and ViTs have opened a new chapter. It may be that they replace ConvNets in the long term.
\end{iband}
{\fontsize{9.4}{12.6}\selectfont\color{fg2}You will most likely meet the architecture in \textbf{large-scale generative image models} --- chapter 17. For small-scale classification, \hilite{ConvNets remain the better bet}\par}\vskip 4pt
\end{frame}

\framekicker{Summary}
\begin{frame}[shrink=25]{What has to survive this chapter}
\begin{isteps}
\item \textbf{Architecture is a choice of hypothesis space} --- it encodes prior knowledge, and bad choices cannot be trained away.
\item \textbf{Modularity, hierarchy, reuse.} Repeated blocks, and a pyramid of growing depth against shrinking size.
\item \textbf{Residual connections} defeat vanishing gradients; project with a 1$\times$1 convolution when shapes differ.
\item \textbf{Batch normalisation} after convolutions, with \inlinecode{use\_bias=False} --- and leave it frozen when fine-tuning.
\item \textbf{Depthwise separable convolutions} encode a stronger prior with \textasciitilde{}8$\times$ fewer parameters, even though the GPU will not run them faster.
\end{isteps}
\end{frame}

\framekicker{Summary}
\begin{frame}[shrink=25]{\ldots{}and the number worth remembering}
\begin{iband}[signal]
Assembling those four practices into a mini Xception took the same dogs-vs-cats problem from \textbf{83.9\% to 90.8\%} while \hilite{halving the parameter count} --- before any hyperparameter tuning.
\end{iband}
\vskip 4pt
\reslink{NOTEBOOK}{03\_mini\_xception.ipynb}{../../course-slides/notebooks/ch09/03\_mini\_xception.ipynb}
\reslink{NEXT}{Chapter 10 --- Interpreting what ConvNets learn}{../ch10/index.html}
\vskip 2pt
\end{frame}

\end{document}
